\documentclass[12pt,a4paper]{article}

\usepackage[utf8]{inputenc}
\usepackage[T1]{fontenc}
\usepackage[spanish,es-nodecimaldot]{babel}
\usepackage{lmodern}
\usepackage{textcomp}
\usepackage{amsmath,amssymb}
\usepackage{booktabs}
\usepackage{array}
\usepackage{float}
\usepackage{graphicx}
\usepackage{xcolor}
\usepackage{microtype}
\usepackage{tikz}
\usepackage[hidelinks]{hyperref}
\usepackage[a4paper,margin=2.5cm]{geometry}

\usetikzlibrary{arrows.meta,babel,calc}

\setlength{\parindent}{0pt}
\setlength{\parskip}{0.65em}
\setcounter{tocdepth}{3}
\renewcommand{\arraystretch}{1.25}

\let\oldsection\section
\renewcommand{\section}{\clearpage\oldsection}

\begin{document}
\hypersetup{pageanchor=false}

\begin{titlepage}
\centering
\textbf{UNIVERSIDAD NACIONAL DE CÓRDOBA}\par
\vspace{0.5em}
\textbf{Facultad de Ciencias Exactas, Físicas y Naturales}\par
\vspace{0.5em}
\textbf{Ingeniería en Computación}\par
\vfill
{\LARGE\bfseries Trabajo Práctico N.\textsuperscript{o} 5\par}
\vspace{1.5em}
{\Large\bfseries Transmisión de Datos: de ICMP y ARP a Sockets TCP entre Dos Computadoras\par}
\vfill
\begin{tabular}{@{}ll@{}}
\textbf{Asignatura:} & Redes de Computadoras \\
\textbf{Grupo:} & MiLANesas \\
\textbf{Alumnos:} &
  \begin{tabular}[t]{@{}ll@{}}
  Tomas Brigido & Kevin Cufre \\
  Rocco Delgado & Facundo Lugones \\
  Javier Siñuka & Agustin Tosolini
  \end{tabular} \\
\textbf{Fecha:} & Septiembre de 2026
\end{tabular}
\vfill
\end{titlepage}

\hypersetup{pageanchor=true}
\pagenumbering{roman}
\tableofcontents
\clearpage
\pagenumbering{arabic}

\section{Ítem 1: ICMP y primer contacto con Wireshark}

\section{Ítem 2: ARP: de una IP a una dirección MAC}

\section{Ítem 3: TCP y UDP ``a mano'' con ncat}

\section{Ítem 4: Servidor TCP mínimo - TP5 Redes de Computadoras}

\end{document}
